
%(BEGIN_QUESTION)
% Copyright 2006, Tony R. Kuphaldt, released under the Creative Commons Attribution License (v 1.0)
% This means you may do almost anything with this work of mine, so long as you give me proper credit

Et krevende spørsmål som noen ganger dukker opp i intervjuer, går omtrent slik:

\vskip 10pt {\narrower \noindent \baselineskip5pt

{\it «Fortell om en hendelse på jobb der du gjorde en feil, og beskriv også hva du gjorde for å rette den.»}

\par} \vskip 10pt

En vanlig feil mange uerfarne kandidater gjør, er å ikke vise til en faktisk hendelse fra eget liv når de svarer på slike spørsmål. I stedet svarer de hypotetisk, og forteller intervjueren hva de {\it kunne} ha gjort {\it hvis} noe slikt {\it skulle skje}.

Forklar hvorfor det er viktig å svare på slike spørsmål med {\it virkelige erfaringer} og ikke hypotetisk. Forklar også konkret hvordan du selv ville svart på dette spørsmålet om feil, og hvilke positive egenskaper ved deg som kommer frem i svaret ditt.

\vskip 50pt

\underbar{file i00745}
%(END_QUESTION)





%(BEGIN_ANSWER)

Dette kan du diskutere med medstudentene dine!

%(END_ANSWER)





%(BEGIN_NOTES)

Igjen er ærlighet beste strategi: ikke lag et oppdiktet scenario for å tilfredsstille intervjueren. Fortell heller om en ekte erfaring som viser evne til å korrigere deg selv og bli bedre. Historien bør ende positivt, slik at den viser vilje og kapasitet for faglig utvikling.

%INDEX% Karriere, intervju: spørsmål om tidligere feil

%(END_NOTES)


